\scanpage{465}\mspage{377}
% archive sheet number: 464
En particulier, on trouve

\uline{Corollaire} 4.14.6. Soient $M,N$ deux catégories accessibles, supposons $M$ stable par petites $\varinjlim$ (i.e $M$ est une cat. \unsure{paratopos}). Alors la catégorie $\underline{\mathrm{Hom}}_{\deleted{\unsure{\Phi}}!}(M,N)$ des foncteurs $f:M\longrightarrow N$ qui commutent aux \inserted{petites} $\varinjlim$ est accessible. Si $N$ est elle aussi un paratopos, $\underline{\mathrm{Hom}}_!(M,N)$ est un paratopos.

\uline{Dém} La première assertion est le cas particulier de 4.14.5 relatif au cas où $\Phi=\Phi_0$, $\Psi=\varnothing$. La deuxième assertion en résulte, puisque $\underline{\mathrm{Hom}}_!(M,N)$ est \unsure{visiblemt|visiblement} stable par petites $\varinjlim$, \unsure{donc|car} c'est \unsure{un|une} paratopos étant accessible (4.8.9 page 207).

\uline{Corollaire} 4.14.7. Soient $\Phi,\Psi,$\inserted{$M$} comme dans 4.14.5, \deleted{et soit $M$ une catégorie accessible telle que $M^\circ$ (catégorie opposée à $M$) soit stable par $\varinjlim$ de type $\Phi$, par $\varprojlim$ de type $\Psi$, i.e. telle que} \inserted{Alors pour \deleted{\unsure{toute}} catégorie accessible $N$, la catégorie}
